% Abstract

\pdfbookmark[1]{Resumen}{Abstract} % Bookmark name visible in a PDF viewer

\begingroup
\let\clearpage\relax
\let\cleardoublepage\relax
\let\cleardoublepage\relax




\chapter*{Resumen} % Abstract name

Este trabajo presenta el diseño, implementación y verificación de \ac{VICON}, un sistema de adquisición y transmisión de imagen en tiempo 
real basado en \ac{FPGA}. El sistema integra el sensor de imagen CMOS MT9V111 con la plataforma Basys~3 (Xilinx Artix-7) y transmite el flujo de vídeo hacia un PC mediante 
el chip FT232H de FTDI operando en modo Synchronous FIFO a 60~MHz.
El diseño Hardware, desarrollado íntegramente en \ac{VHDL}, comprende un controlador I2C para la configuración del sensor, un módulo de captura de trama con gestión 
de marcadores de protocolo, un controlador de alta velocidad para el FT232H con arquitectura de doble proceso para TX y RX y un procesador de comandos con 
sincronización entre dominios de reloj (\ac{CDC}). La verificación se realiza mediante simulación con \ac{UVVM}, 
incluyendo agentes de simulación para el sensor I2C y el chip FT232H. El software de visualización, desarrollado en C++ con Qt y OpenCV, recibe el flujo de imagen y permite enviar comandos de control a la \ac{FPGA} en tiempo real.
El sistema completo ha sido validado sobre hardware real, demostrando la transmisión de imagen en escala de grises a resolución del estándar \ac{VGA} (640$\times$480) y la correcta ejecución de todos los comandos de control.
\endgroup
\providecommand{\keywords}[1]
{
	\small	
	\textbf{\textit{Keywords---}} #1
}			

\vfill